% =====================================================================
%  Section 3.5 — Proposed Improvements and Expected Gains
%  Insert after Section 3.4 (Limitation).
%
%  NOTE: every number in this section is a PROJECTION derived from the
%  failure modes observed in Section 3.3, not a measured result.
%  The wording below states this explicitly and should not be softened.
% =====================================================================

\section{Proposed Improvements and Expected Gains}
\label{sec:expected-gains}

The evaluation in Section~\ref{sec:qualitative-eval} identified three recurring
failure modes: weak handling of multi-condition prompts, insufficient factual
grounding, and inconsistent enforcement of negative constraints. This section
maps each failure mode onto a concrete architectural change and estimates the
gain that change is expected to produce under the same evaluation protocol.

\paragraph{Implementation status.}
Five of the changes below have since been implemented, and their effect on
retrieval and ranking has been measured directly; those results are reported in
Section~\ref{sec:retrieval-eval} and reconciled against the projections in
Section~\ref{sec:projection-vs-measurement}. The judge scores projected in this
section have \emph{not} been re-measured, because doing so requires re-running
the protocol of Section~\ref{sec:qualitative-eval} on the revised system.
Everything in Tables~\ref{tab:expected-gains} and~\ref{tab:expected-dimensions}
therefore remains a projection.

\paragraph{Status of the figures reported here.}
All scores in this section are \emph{projections}, not measurements. They are
derived analytically from the failure modes documented in
Section~\ref{sec:qualitative-eval} by asking, for each observed error, whether
the proposed mechanism would have prevented it. They have not been obtained by
re-running the evaluation, and they are reported here only to motivate the
prioritisation of future work. We deliberately distinguish between changes whose
effect is \emph{deterministic} (the mechanism structurally forbids the observed
error) and changes whose effect is \emph{probabilistic} (the mechanism makes the
error less likely but cannot exclude it), and we assign confidence levels
accordingly.

\subsection{Proposed Changes}
\label{sec:proposed-changes}

\begin{description}

  \item[C1 --- Hybrid retrieval with rank fusion.] \hfill\textsc{implemented; confirmed}\\
  The current pipeline treats lexical (BM25) retrieval as a fallback that is
  reached only when vector retrieval is unavailable, so in practice the lexical
  index is never queried. Both retrievers would instead run on every query and
  their ranked outputs would be merged using Reciprocal Rank Fusion (RRF), which
  combines documents by rank rather than by score and therefore requires no
  normalisation between the two incomparable score scales. The motivation is
  that the two retrievers fail in complementary ways: dense retrieval degrades
  on proper nouns, where a title such as \emph{Stardew Valley} is embedded close
  to its genre neighbours, whereas lexical retrieval degrades on paraphrase and
  mood-based queries.

  \item[C2 --- Cross-encoder re-ranking.] \hfill\textsc{not implemented}\\
  Section~\ref{sec:ranking} currently orders candidates by their retrieval score
  and truncates, which is not re-ranking in a meaningful sense. Generative
  re-ranking with the local LLM was implemented but disabled because it added
  approximately 550\,s per query on CPU. The cost is dominated by autoregressive
  decoding over a long prompt rather than by the ranking task itself. A
  cross-encoder re-ranker scores each query--candidate pair in a single forward
  pass and produces no text, reducing the same operation to a few seconds on the
  same hardware. Because the cross-encoder attends to the full query against
  each candidate, it is specifically strong on multi-clause queries, where a
  single dense vector tends to be dominated by the most salient term.

  \item[C3 --- Structured generation with identifier validation.] \hfill\textsc{not implemented}\\
  The answer-generation step currently emits free text from which titles are
  recovered by regular expressions and matched back to records by name prefix.
  Constrained decoding against a JSON schema would instead force the model to
  return application identifiers, which are then validated against the candidate
  set before display. Any identifier outside the candidate set is discarded and
  counted.

  \item[C4 --- Constraint decomposition and pre-filtering.] \hfill\textsc{not implemented}\\
  Hard constraints (price, platform, multiplayer support) and soft preferences
  (mood, theme) are currently handled by the same substring-matching module, and
  a hard constraint is silently dropped whenever it would empty the candidate
  list. Decomposing the query into a structured filter and a semantic component
  --- either by LLM tool calling or by the existing rule-based parser as a
  fallback --- allows hard constraints to be pushed down into the SQL query, so
  that filtering occurs \emph{before} retrieval rather than after it.

  \item[C5 --- Review-aware quality scoring.] \hfill\textsc{implemented; confirmed}\\
  The current popularity term uses the positive review count alone and therefore
  ignores negative reviews entirely. Replacing it with the Wilson lower bound of
  the positive ratio, blended with log review volume, penalises titles whose
  high ratings rest on very few observations while still rewarding
  well-established games.

  \item[C6 --- Diversity constraints.] \hfill\textsc{implemented; not confirmed}\\
  Caps on the number of results sharing a developer or a primary tag prevent a
  shortlist from collapsing onto near-identical titles. The constraint is
  defined so that it may only reorder results, never reduce their number.

  \item[C7 --- Index corpus rebuild.] \hfill\textsc{implemented; partly confirmed}\\
  The dense index covered 18,449 of 39,175 titles while the lexical index
  covered all of them, and it applied neither of the embedding model's required
  task prefixes. Rebuilding raised coverage to 30,693 --- excluding only titles
  whose rating metadata is absent, for which the review-aware signal cannot be
  computed --- added the prefixes, and reordered review sampling by community
  votes rather than length.

  \item[C8 --- Reviews in the lexical corpus.] \hfill\textsc{implemented; confirmed}\\
  The dense documents carried player reviews and the lexical corpus did not,
  so queries phrased as descriptions of play had almost nothing to match
  lexically. Both corpora now draw review text through the same helper.

\end{description}

\subsection{Projected Results by Prompt Category}
\label{sec:projected-results}

Table~\ref{tab:expected-gains} maps each prompt category from
Table~\ref{tab:test-prompts} onto the changes expected to affect it.

\begin{table}[h]
\centering
\caption{Projected scores by prompt category under the evaluation protocol of
Section~\ref{sec:qualitative-eval}. \textbf{All ``Projected'' values are
estimates derived from observed failure modes, not measurements.}}
\label{tab:expected-gains}
\small
\begin{tabular}{llccclc}
\hline
\textbf{\#} & \textbf{Category} & \textbf{Current} & \textbf{Projected} &
\textbf{$\Delta$} & \textbf{Primary driver} & \textbf{Conf.} \\
\hline
1 & Simple / Vague       & 88 & 92 & $+4$  & C5                & High   \\
2 & Detailed Descriptive & 50 & 75 & $+25$ & C2, C4            & Medium \\
3 & Negation-Constrained & 61 & 78 & $+17$ & C4, C2            & Medium \\
4 & Hard-Constrained     & 50 & 87 & $+37$ & C3, C4            & High   \\
5 & Analogy              & 52 & 73 & $+21$ & C1, C2            & Low    \\
\hline
  & \textbf{Mean}        & \textbf{60.2} & \textbf{81.0} & $\mathbf{+20.8}$ & & \\
\hline
  & ChatGPT reference    & 95.4 & 95.4 & --- & & \\
  & Gap to reference     & $-35.2$ & $-14.4$ & --- & & \\
\hline
\end{tabular}
\end{table}

\paragraph{Category 4 (Hard-Constrained), $50 \rightarrow 87$, high confidence.}
This category shows the largest projected gain because its dominant error is
deterministic rather than a matter of degree. Two of the five recommended
titles, \emph{Wolvesville} and \emph{Sonder}, could not be verified on Steam,
which forfeits most of the 30-point verifiability dimension. Under C3 these
titles cannot be returned at all: they are absent from the underlying snapshot,
therefore absent from the candidate set, therefore rejected by identifier
validation. The remaining loss on the relevance dimension is addressed by C4,
which enforces ``free'' and ``co-operative'' as SQL predicates before retrieval
rather than as post-hoc filters that are dropped when they would empty the
candidate list. We expect the verifiability dimension to approach its maximum.

\paragraph{Category 2 (Detailed Descriptive), $50 \rightarrow 75$, medium confidence.}
The observed failure is the loss of one clause of a two-clause query: the system
matched ``intellectually stimulating'' while ignoring ``require some gaming
experience''. This is the canonical case in which cross-encoders outperform
bi-encoders (C2), since the re-ranker evaluates the complete query against each
candidate and can penalise candidates satisfying only one clause, whereas a
single query embedding is dominated by its most salient term. C4 reinforces this
by representing the clauses as separate criteria. The projection is moderate
rather than large because clause coverage is improved but not guaranteed.

\paragraph{Category 3 (Negation-Constrained), $61 \rightarrow 78$, medium confidence.}
Two distinct errors occurred. The recommendation of a single-player title under a
``play with friends'' requirement is a hard-constraint violation and is addressed
deterministically by C4, using the multiplayer category field as a pre-filter.
The recommendation of a visually cute title under a ``no cozy or cute'' exclusion
is a semantic negation, which is addressed only probabilistically. Dense
embeddings are known to be weak at negation, since a phrase and its negation
share almost all of their lexical and topical content and therefore lie close in
embedding space; the cross-encoder and an explicit exclusion field improve on
this but cannot eliminate it. The projection is therefore deliberately more
conservative than for Categories 2 and 4.

\paragraph{Category 5 (Analogy), $52 \rightarrow 73$, low confidence.}
The system anchored on ``more challenging'' and lost the ``like Stardew Valley''
reference, returning titles such as \emph{Diablo IV} and \emph{ICARUS}. Treating
an explicit title reference as an entity anchor --- resolving it against the
catalogue and using its tag profile as the semantic query, with the remaining
clause as a modifier --- is expected to recover much of this, and C1 helps
because exact matching of a proper noun is precisely where lexical retrieval
outperforms dense retrieval. This projection carries the lowest confidence
because analogical reasoning depends primarily on world knowledge about how
games actually play, which is the dimension on which a small local model is least
competitive.

\paragraph{Category 1 (Simple / Vague), $88 \rightarrow 92$, high confidence but limited headroom.}
Performance here is already close to the ceiling. The residual gap is
concentrated in the projected-satisfaction dimension: the recommended titles are
valid but obscure, whereas the reference system returned widely acclaimed ones.
C5 addresses exactly this by promoting, among comparably relevant candidates,
those whose quality is supported by a large volume of reviews.

\subsection{Projected Effect by Scoring Dimension}
\label{sec:projected-dimensions}

\begin{table}[h]
\centering
\caption{Projected effect by scoring dimension, averaged across the five
prompts. \textbf{Estimates, not measurements.}}
\label{tab:expected-dimensions}
\small
\begin{tabular}{lccl}
\hline
\textbf{Dimension} & \textbf{Current (est.)} & \textbf{Projected} & \textbf{Basis} \\
\hline
Verifiability (30)          & $\approx 22$ & $\approx 29$ & C3: structural, deterministic \\
Relevance (30)              & $\approx 17$ & $\approx 24$ & C2, C4: probabilistic \\
Projected satisfaction (40) & $\approx 21$ & $\approx 28$ & C5: bounded by model judgement \\
\hline
\end{tabular}
\end{table}

The dimension-level view clarifies where the gains originate. The verifiability
dimension is the only one for which a structural guarantee is available:
identifier validation converts hallucination from a probabilistic failure into
an impossible one. Gains on relevance and projected satisfaction rest on
improved ranking and are therefore expected but not guaranteed.

\subsection{Projections Against Measurements}
\label{sec:projection-vs-measurement}

The retrieval-level evaluation in Section~\ref{sec:retrieval-eval} tests five of
these changes on metrics the projections above do not use, so the two cannot be
compared numerically. What can be compared is the \emph{ordering}: which changes
were expected to matter most, and which turned out to.

\begin{table}[h]
\centering
\caption{Expected against measured standing for the implemented changes.
Measurements are retrieval-level; the projected judge scores remain unmeasured.}
\label{tab:projection-vs-measurement}
\small
\begin{tabular}{llll}
\hline
\textbf{Change} & \textbf{Expected} & \textbf{Measured} & \textbf{Agreement} \\
\hline
C1 Rank fusion        & moderate, analogy category   & confirmed, $3.5\sigma$ on known-item & yes \\
C5 Quality scoring    & small, vague category only   & confirmed, $2.9$--$3.6\sigma$ on topical & understated \\
C6 Diversity          & improves shortlist usability & not confirmed; binds on 14\% of queries & overstated \\
C7 Prefixes           & certain gain, low risk       & no measurable effect on either track & overstated \\
C7 Coverage           & not separately projected     & the whole of C7's measured gain & understated \\
C8 Lexical reviews    & not foreseen                 & largest single effect: $0.000 \rightarrow 0.133$ & missed \\
\hline
\end{tabular}
\end{table}

\paragraph{Three of six projections were wrong, in instructive ways.}
The embedding prefix was described above as a certain, low-risk repair; it
produces no measurable gain on either evaluation track. The diversity
constraint was expected to improve shortlist usability; it alters the top five
on 7 of 50 queries and leaves intra-list diversity unchanged, because fused
shortlists are already varied enough that its caps rarely bind. Conversely,
review-aware quality scoring was projected to matter only for vague prompts and
is in fact the strongest confirmed component on topical queries.

\paragraph{The largest measured effect was not projected at all.}
That the lexical corpus omitted review text was not identified when this section
was written. It surfaced only because the lexical configuration scored exactly
0.000 on every known-item metric, which is an implausible result for a working
retriever and therefore prompted a control test. Repairing it moved Recall@50
from 0.000 to 0.133 --- larger than any change that had been projected.

\paragraph{Consequence for the projections that remain.}
C2, C3 and C4 are still unimplemented and their projected judge scores stand
unverified. Given that half of the projections for the implemented changes were
wrong in direction or magnitude, those figures should be read as a
prioritisation aid rather than as forecasts. The mechanism-level reasoning
behind them --- particularly the distinction between deterministic and
probabilistic fixes --- survives better than the numbers did: the one change
whose effect was argued to be structural rather than statistical, C3's
identifier validation, has not yet been tested.

\subsection{Expected Residual Gap}
\label{sec:residual-gap}

The projections close approximately 59\% of the measured gap to the commercial
reference system but do not eliminate it, and we do not expect them to. As noted
in Section~\ref{sec:eval-summary}, the reference system benefits from live web
access, a substantially larger model, and world knowledge about how individual
games actually play. The changes proposed here address retrieval coverage,
ranking quality, and factual grounding; they do not confer knowledge of what
makes a game satisfying to play. We therefore expect a persistent residual gap
concentrated in the projected-satisfaction dimension and in the analogy
category, both of which depend on that knowledge rather than on pipeline
design.

This bound is itself a useful finding: it separates the portion of the observed
gap that is attributable to system architecture, and is therefore addressable
within the constraints of this assignment, from the portion attributable to
model scale and data freshness, which is not.

\subsection{Proposed Revisions to the Evaluation Protocol}
\label{sec:eval-protocol-revision}

Verifying the projections above requires a stronger protocol than the one used
in Section~\ref{sec:qualitative-eval}, for three reasons.

\begin{enumerate}
  \item \textbf{Sample size.} Five prompts, one per category, give a mean with a
  wide confidence interval; a single prompt swing of ten points moves the
  reported mean by two. We propose four to five prompts per category, for a
  total of twenty to twenty-five.

  \item \textbf{Judge variance.} An LLM judge is not deterministic, so a single
  scoring pass conflates system differences with judge noise. We propose running
  the judge three times per comparison and reporting the mean and standard
  deviation, so that a claimed improvement can be checked against the
  measurement noise.

  \item \textbf{Diagnostic resolution.} A single composite score indicates that
  a result is poor but not which stage caused it. We propose adding automatic
  retrieval metrics --- Recall@50 for the retrieval stage and nDCG@10 for the
  ranking stage --- computed against relevance labels derived from tag overlap
  and validated against a manually annotated subset. Separating the two makes it
  possible to distinguish a candidate that was never retrieved from one that was
  retrieved and ranked too low, which the current protocol cannot do.
\end{enumerate}

Together these changes would allow the projections in
Table~\ref{tab:expected-gains} to be confirmed or refuted per component, rather
than assessed only in aggregate.

\paragraph{What has since been done.}
The third proposal --- separating retrieval from ranking with automatic metrics
--- has been implemented and is reported in Section~\ref{sec:retrieval-eval},
with 550 queries across two independent sources of ground truth rather than the
20 to 25 suggested here. The first two, expanding the judge prompt set and
scoring each comparison several times, remain outstanding; they apply to the
judge protocol, which has not been re-run on the revised system.
